\section{Proof Templates, Oracles, And Handles}

Before we define fields, tensors, or codes, we need to fix one vocabulary point.
The word ``oracle'' is used in several nearby but different proof templates.  In
one place it means an ideal proof string that the verifier may open by address.
In another place it means a large input that the verifier may open by address.
In a multilinear protocol it means an object that evaluates a multilinear
polynomial at field points.  After cryptographic compilation it may mean a
Merkle message-oracle handle that authenticates coordinate openings.

Those are not the same interface.  This section gives each one a name.  Later
sections may say that an object is virtual, committed, encoded, or holographic,
but it must still implement one of the interfaces below.  This is the discipline
that keeps the Blaze compiler stack from becoming a soup of unnamed oracles.

In the idealized proof templates in this section, an oracle is best viewed as a
trusted functionality, or an extra abstract party.  The prover, indexer, or
environment commits a fixed object to that functionality.  The functionality
returns a handle.  Later, the verifier uses the handle to ask the functionality
for selected information.  The functionality is trusted to answer consistently
from the committed object, but it is not trusted to make that object correct.
Correctness is still enforced only by the verifier's checks.  In a real
protocol, these ideal functionalities are instantiated by concrete mechanisms
such as Merkle commitments, polynomial commitments, or transcript-bound virtual
opening rules.

The source path is: IOPs from BCS16~[BCS16], IOPPs from RVW13~[RVW13],
holographic IOPs from Fractal~[COS20], multilinear IOPs and encoded-MLE
compilation from Blaze~[Blaze24], polynomial-oracle terminology and batching from
HyperPlonk~[CBBZ23], and Fiat--Shamir/RBR soundness from
CMS19/CCH19~[CMS19,CCH19].

\subsection{Relations And Roles}

The basic object proved by a proof system is a relation.  Following the
relation/language convention used in BCS16~[BCS16], let
\[
  R\subseteq \mathcal X\times\mathcal W
\]
be a set of valid instance-witness pairs.  Its language is
\[
  L(R)
  =
  \{x\in\mathcal X:\exists w\in\mathcal W\text{ such that }(x,w)\in R\}.
\]
The verifier receives an instance \(x\).  The honest prover additionally knows
a witness \(w\).  Completeness says that valid pairs are accepted.  Soundness
says that if \(x\notin L(R)\), no malicious prover should make the verifier
accept except with small probability.

The statement being proved need not be a short string.  It may contain one or
more large strings.  When the verifier receives such a large string by handle,
rather than reading it in full, we will call that handle an \emph{input oracle}
for the protocol.  Logically, the input is the underlying string
\(Y\in\Sigma^D\).  Operationally, the verifier only sees a handle \(O_Y\) and
opens a sparse set of coordinates
\[
  \mathsf{Open}(O_Y,a)\to Y[a]\quad\text{or}\quad\bot.
\]
This sparse-access treatment of a statement component is the holographic idea.
The word describes the access pattern, not the semantic role of \(Y\).

The handle may be fixed before the protocol starts, or it may be produced by an
earlier phase and then used as an input oracle by a later phase.  This matters
for Blaze/BaseFold: a Blaze phase can create a committed or virtual oracle, and
a later encoded-MLE phase can treat that existing oracle as its input string.
The later verifier's logical input is still the underlying string \(Y\), even
though it only queries a few locations through \(O_Y\).

One example is the standard holographic IOP use case.  Following the
holographic-IOP convention used by Fractal~[COS20], an indexed relation has
triples
\[
  R\subseteq \mathcal I\times\mathcal X\times\mathcal W.
\]
Here \(i\in\mathcal I\) is the index, \(x\in\mathcal X\) is the online
instance, and \(w\in\mathcal W\) is the witness.  The index may be too large for
the verifier to read directly, so a holographic protocol gives the verifier
opening access to an encoding of \(i\).

Another example is the Blaze/BaseFold setting.  Blaze also uses holographic
access for the long string being proved about.  The paper calls this string the
implicit input~[Blaze24].  In this explainer we will name the string locally
instead of relying on that phrase.  The same coordinate-handle interface can
therefore serve several roles.

\begin{description}[leftmargin=0.32in,style=nextline]
  \item[Index.]
        A structured public parameter \(i\), such as a circuit, code
        description, or permutation table.  The verifier may open coordinates
        of an encoding \(\mathsf{idx}(i)\).
  \item[Explicit input.]
        The verifier-held short input \(x\), read directly.
  \item[Long statement string.]
        A named string \(Y\in\Sigma^D\) that is part of the statement and is
        fixed behind a coordinate handle.  A later section must name the
        concrete string, its domain \(D\), and the relation or proximity
        condition being checked against it.
  \item[Prover oracle string.]
        A string sent by the prover during an IOP or IOPP round.  It also gets
        a coordinate handle, but it is proof material rather than statement
        data.
  \item[Witness.]
        A prover-held value \(w\) proving that the named statement strings and
        explicit input satisfy the relation.
\end{description}

This section intentionally stops at the role level.  Later sections instantiate
these roles with concrete tensors, domains, codes, and evaluation claims before
using them in protocol statements.

\subsection{From IPs To IOPs}

An interactive proof (IP) is ordinary interaction, as recalled in BCS16's route
to IOPs~[BCS16].  The verifier sends and receives messages, reads every
received message in full, and eventually outputs accept or reject.  In a
public-coin IP, every verifier message is sampled from a declared public
randomness domain.

A probabilistically checkable proof (PCP) changes the access pattern~[BCS16].
The proof is a fixed string, and the verifier opens only selected coordinates of
that string.  There is no later prover adaptation.

An interactive oracle proof (IOP), introduced by BCS16~[BCS16], combines those
two ideas.  The prover fixes message strings across rounds, and the verifier
later opens selected coordinates.  We model this access through the ideal
message-oracle functionality in Figure~\ref{fig:func-msg-oracle}.  Operationally,
the functionality binds a fixed string and answers coordinate openings.  BaseFold
and Blaze use the public-coin, prover-first shape: the prover first fixes a
string, the verifier samples public challenges, and the prover fixes more
strings~[BaseFold23,Blaze24].

\begin{protocolbox}{Functionality: Message Oracle}{fig:func-msg-oracle}
Let \(\Sigma\) be an alphabet and \(D\) a finite address set.

\smallskip
\noindent\(\underline{\mathcal F_{\mathrm{MsgOracle}}.\mathsf{Commit}(\mathsf{source}:\pi\in\Sigma^D,\ \verifier:\bot)\to O_\pi}\)

\smallskip
\begin{enumerate}[leftmargin=0.24in,label=\textbf{M\arabic*.},itemsep=0.08em,topsep=0.1em,parsep=0pt]
  \item \(\mathcal F_{\mathrm{MsgOracle}}\): store the input \(\pi\)
        under a fresh handle \(O_\pi\).
  \item \(\mathcal F_{\mathrm{MsgOracle}}\): return \(O_\pi\).
\end{enumerate}

\smallskip
\noindent\(\underline{\mathcal F_{\mathrm{MsgOracle}}.\mathsf{Open}(\verifier:(O_\pi,a\in D))\to(\verifier:o\in\Sigma\cup\{\bot\})}\)

\smallskip
\begin{enumerate}[leftmargin=0.24in,label=\textbf{O\arabic*.},itemsep=0.08em,topsep=0.1em,parsep=0pt]
  \item \(\mathcal F_{\mathrm{MsgOracle}}\): if \(O_\pi\) is a stored handle
        for \(\pi\in\Sigma^D\), set \(o\gets\pi[a]\); otherwise set
        \(o\gets\bot\).
  \item \(\mathcal F_{\mathrm{MsgOracle}}\): return \((\verifier:o)\).
\end{enumerate}
\end{protocolbox}

\begin{protocolbox}{Template: Public-Coin IOP}{fig:template-public-coin-iop}
Let \(R\subseteq\mathcal X\times\mathcal W\) and let
\(\varepsilon\in[0,1]\) be the declared soundness error.  There is no external
openable input object in this template; every message-oracle handle below is
created by the prover during the protocol.

\smallskip
\noindent\(\underline{\mathsf{PublicCoinIOP}[R,\varepsilon](\verifier:x\in\mathcal X,\ \prover:(x,w))\to\{0,1\}}\)

\smallskip
\begin{enumerate}[leftmargin=0.24in,label=\textbf{I\arabic*.},itemsep=0.08em,topsep=0.1em,parsep=0pt]
  \item \(\prover\): compute
        \(\pi_0=f_0(x,w)\in\Sigma_0^{D_0}\).
  \item \(\prover,\verifier,\mathcal F_{\mathrm{MsgOracle}}\): run
        \(\mathcal F_{\mathrm{MsgOracle}}.\mathsf{Commit}(\pi_0)\),
        producing \(O_0\).
  \item For each round \(j\in[T]\):
        \begin{enumerate}[leftmargin=0.2in,label=\textbf{I\arabic{enumi}.\arabic*.},itemsep=0.05em,topsep=0.05em,parsep=0pt]
          \item \(\verifier\): sample public challenge \(\alpha_j\) from its
                declared challenge domain.
          \item \(\prover\): compute
                \(\pi_j=f_j(x,w,\alpha_1,\ldots,\alpha_j)\in\Sigma_j^{D_j}\).
          \item \(\prover,\verifier,\mathcal F_{\mathrm{MsgOracle}}\): run
                \(\mathcal F_{\mathrm{MsgOracle}}.\mathsf{Commit}(\pi_j)\),
                producing \(O_j\).
        \end{enumerate}
  \item \(\verifier\): choose finite opening sets \(Q_j\subseteq D_j\) for
        \(j\in\{0,\ldots,T\}\).
  \item \(\verifier\): for each \(j\in\{0,\ldots,T\}\) and each
        \(a\in Q_j\), call
        \(\mathcal F_{\mathrm{MsgOracle}}.\mathsf{Open}(O_j,a)\).
  \item \(\verifier\): output \(0\) or \(1\) from the transcript, public
        challenges, and opened values.
  \item \(\mathsf{Completeness}\): if \((x,w)\in R\), the honest prover is
        accepted.
  \item \(\mathsf{Soundness}\): if there is no \(w\in\mathcal W\) such that
        \((x,w)\in R\), every malicious prover is accepted with probability at
        most \(\varepsilon\).
\end{enumerate}
\end{protocolbox}

Figure~\ref{fig:func-msg-oracle} is an ideal message oracle for a fixed
openable string.  The committing source may be the prover, an indexer, the
environment, or a prior protocol.  Operationally, the functionality is
commitment-like: it stores the whole string and later answers coordinate
openings.
The protocol using a handle must specify when the handle is fixed.  The same
open interface can also be implemented by a Merkle message-oracle handle, an
indexer-provided handle, or a virtual handle.

The base IOP template in Figure~\ref{fig:template-public-coin-iop} is an exact
membership template: the verifier is trying to decide whether the explicit input
\(x\) has some witness \(w\) in the relation \(R\).  The next two templates add
external handles, but they use those handles for different semantic jobs.

\subsection{IOPPs And Proximity}

An interactive oracle proof of proximity (IOPP) adds one more openable object:
the verifier has oracle access to a large input string and wants to know whether
that string is close to the valid set.  This is the proximity setting of
RVW13~[RVW13] and the codeword-proximity setting used by BaseFold~[BaseFold23].
The IOPP template takes this large input as an external message-oracle handle
\(O_y\).  The handle must implement the open interface of
Figure~\ref{fig:func-msg-oracle}; the template does not care whether it was
produced by \(\mathcal F_{\mathrm{MsgOracle}}.\mathsf{Commit}\), by an indexer,
by a Merkle message-oracle handle, or by a virtual construction.

The semantic difference from an ordinary IOP is that \(y\) itself is the object
being tested.  The verifier is not promised exact validity or exact invalidity.
Instead, the template has a gap: valid \(y\)'s should accept, far \(y\)'s should
reject, and close-but-invalid \(y\)'s are outside the promise.

For two strings \(y,y'\in\Sigma^D\), their relative Hamming distance is
\[
  \Delta(y,y')
  =
  \frac{|\{a\in D:y[a]\ne y'[a]\}|}{|D|}.
\]
For a set \(S\subseteq\Sigma^D\), define
\[
  \Delta(y,S)=\min_{s\in S}\Delta(y,s).
\]

\begin{protocolbox}{Ideal Functionality: Proximity To A Language}{fig:func-prox-language}
\noindent Public parameters: a finite domain \(D\), an alphabet \(\Sigma\), a
language \(L\subseteq\Sigma^D\), and a proximity threshold
\(\delta\in[0,1]\).

\smallskip
\noindent\(\underline{\mathcal F_{\mathrm{Prox}}^{L,\delta}.\mathsf{Commit}(\mathsf{source}:y\in\Sigma^D,\ \verifier:\bot)\to I}\)

\smallskip
\begin{enumerate}[leftmargin=0.24in,label=\textbf{S\arabic*.},itemsep=0.08em,topsep=0.1em,parsep=0pt]
  \item \(\mathcal F_{\mathrm{Prox}}^{L,\delta}\): choose a fresh handle \(I\).
  \item \(\mathcal F_{\mathrm{Prox}}^{L,\delta}\): store the committed string
        \(y\in\Sigma^D\) under \(I\).
  \item \(\mathcal F_{\mathrm{Prox}}^{L,\delta}\): return \(I\).
\end{enumerate}

\smallskip
\noindent\(\underline{\mathcal F_{\mathrm{Prox}}^{L,\delta}.\mathsf{Check}(\verifier:I)\to(\verifier:b\in\{0,1\})}\)

\smallskip
\begin{enumerate}[leftmargin=0.24in,label=\textbf{C\arabic*.},itemsep=0.08em,topsep=0.1em,parsep=0pt]
  \item \(\mathcal F_{\mathrm{Prox}}^{L,\delta}\): retrieve \(y\) stored under \(I\).
  \item \(\mathcal F_{\mathrm{Prox}}^{L,\delta}\): set
        \(\Delta_y\gets\Delta(y,L)\).
  \item \(\mathcal F_{\mathrm{Prox}}^{L,\delta}\): if \(\Delta_y=0\), set
        \(b\gets1\).
  \item \(\mathcal F_{\mathrm{Prox}}^{L,\delta}\): if \(\Delta_y>\delta\), set
        \(b\gets0\).
  \item \(\mathcal F_{\mathrm{Prox}}^{L,\delta}\): if
        \(0<\Delta_y\le\delta\), let the corrupt source, or equivalently the
        simulator for that source, choose \(b\in\{0,1\}\).
  \item \(\mathcal F_{\mathrm{Prox}}^{L,\delta}\): return \((\verifier:b)\).
\end{enumerate}
\end{protocolbox}

Figure~\ref{fig:func-prox-language} is the functionality view of proximity.
An IOPP is a protocol realization of this functionality for a chosen language
\(L\), threshold \(\delta\), and soundness error \(\varepsilon\), using only
sublinear oracle openings of the committed string.  In applications \(L\) may
depend on the public input \(x\); then the instantiated language is
\(L_x=S_{\mathrm{valid}}(x)\).

\begin{protocolbox}{Ideal Functionality: Encoded MLE Opening For A Codeword Handle}{fig:func-encmle-code-opening}
\noindent Public parameters: a field \(\F\), a message dimension \(k\), a
codeword dimension \(n\), a code
\[
  C:\F^{\B^k}\to\F^{\B^n},
\]
and a proximity threshold \(\delta\in[0,1]\).  The handle \(O_y\) implements
\(\mathcal F_{\mathrm{MsgOracle}}.\mathsf{Open}\) for a fixed string
\(y\in\F^{\B^n}\).

\smallskip
\noindent\(\underline{\mathcal F_{\EncMLE[C]}^\delta.\mathsf{Open}(\verifier:(O_y:\B^n\to\F\cup\{\bot\},\ z\in\F^k,\ v\in\F),\ \prover:m\in\F^{\B^k})\to(\verifier:a\in\{0,1\})}\)

\smallskip
\begin{enumerate}[leftmargin=0.24in,label=\textbf{R\arabic*.},itemsep=0.08em,topsep=0.1em,parsep=0pt]
  \item \(\mathcal F_{\EncMLE[C]}^\delta\): read the fixed string
        \(y\in\F^{\B^n}\) behind \(O_y\).
  \item \(\mathcal F_{\EncMLE[C]}^\delta\): compute
        \[
          \Delta_m\gets\Delta(y,C(m)),
          \qquad
          e_m\gets \MLE(m)(z).
        \]
  \item If \(\Delta_m=0\) and \(e_m=v\), return \((\verifier:1)\).
  \item If \(\Delta_m>\delta\) or \(e_m\ne v\), return \((\verifier:0)\).
  \item If \(0<\Delta_m\le\delta\) and \(e_m=v\), let the corrupt source, or
        equivalently the simulator for that source, choose
        \(a\in\{0,1\}\), and return \((\verifier:a)\).
\end{enumerate}
\end{protocolbox}

This is the generic encoded multilinear-evaluation functionality for a code
\(C\).  The verifier's input is not a fresh commitment operation; it is an
already-fixed coordinate oracle \(O_y\).  The functionality asks whether the
string \(y\) behind that oracle is an encoding of a message \(m\), and whether
the multilinear extension of that same \(m\) evaluates to \(v\) at \(z\).
Expanded as an exact relation,
\[
  ((z,v),y;m)\in\EncMLE[C]
  \qquad\Longleftrightarrow\qquad
  y=C(m)\ \text{ and }\ \MLE(m)(z)=v.
\]
The functionality above is the proximity version with \(y\) named by the input
oracle handle \(O_y\).  It does not give the verifier arbitrary MLE access to
\(y\), and it does not say that \(y\) is exactly a codeword.  It is the ideal
specification that a concrete encoded-evaluation IOPP, such as the
Blaze/BaseFold backend opening protocol, attempts to realize.

Terminology in the literature is not uniform.  FRI is usually described as an
IOPP for Reed--Solomon proximity, and a FRI-based PCS adds an evaluation-opening
reduction around that proximity test.  BaseFold describes the same semantic
task as a BaseFold IOPP combined with sumcheck for a multilinear evaluation
claim.  Blaze packages the task as the relation \(\RMLE[C]\).  In this paper,
\(\EncMLE[C]\) is our native name for the common functionality, and
\(\RMLE[C]\) is treated as Blaze's notation for the same relation.

An ordinary MLPCS is immediate once an encoded-MLE opening protocol is
available.  To commit a tensor \(h\in\F^{\B^k}\), the prover forms the encoded
string \(y=C(h)\in\F^{\B^n}\), commits that string to the coordinate-oracle
layer, and receives a handle \(O_y\).  The MLPCS handle is the pair
\((O_y,C)\).  To check an evaluation claim \(\MLE(h)(z)=v\), the verifier runs
\(\mathcal F_{\EncMLE[C]}^\delta.\mathsf{Open}(O_y,z,v)\).  Thus
\(\EncMLE[C]\) is the encoded opening primitive; MLPCS is the usual
commit-then-open interface obtained by placing a commit step in front of that
primitive.

The only subtlety is the proximity/effective-object split.  The literal string
behind \(O_y\) need not equal a codeword in a compiled protocol.  If it is close
to the code, accepted openings are interpreted with respect to the nearby
decoded message \(h^\dagger\).  If it is far, the encoded-MLE opening rejects
except with its stated soundness error.  This is the binding reason that a
prover cannot freely answer different MLPCS openings as different tensors after
the encoded string has been fixed.

\begin{protocolbox}{Template: IOPP}{fig:template-iopp}
Let \(S_{\mathrm{valid}}(x)\subseteq\Sigma^D\) be the valid set for the named
long statement string under explicit input \(x\), let \(\delta\in[0,1]\) be the proximity radius, and
let \(\varepsilon\in[0,1]\) be the declared soundness error.  The handle
\(O_y\) opens a fixed string \(y\in\Sigma^D\).

\smallskip
\noindent\(\underline{\mathsf{IOPP}[S_{\mathrm{valid}},\delta,\varepsilon](\verifier:(x,O_y),\ \prover:(x,w))\to\{0,1\}}\)

\smallskip
\begin{enumerate}[leftmargin=0.24in,label=\textbf{P\arabic*.},itemsep=0.08em,topsep=0.1em,parsep=0pt]
  \item \(O_y\): fixed message-oracle handle implementing
        \(\mathcal F_{\mathrm{MsgOracle}}.\mathsf{Open}\).
  \item \(\prover\): compute
        \(\pi_0=f_0(x,w)\in\Sigma_0^{D_0}\).
  \item \(\prover,\verifier,\mathcal F_{\mathrm{MsgOracle}}\): run
        \(\mathcal F_{\mathrm{MsgOracle}}.\mathsf{Commit}(\pi_0)\),
        producing \(O_0\).
  \item For each round \(j\in[T]\):
        \begin{enumerate}[leftmargin=0.2in,label=\textbf{P\arabic{enumi}.\arabic*.},itemsep=0.05em,topsep=0.05em,parsep=0pt]
          \item \(\verifier\): sample public challenge \(\alpha_j\) from its
                declared challenge domain.
          \item \(\prover\): compute
                \(\pi_j=f_j(x,w,\alpha_1,\ldots,\alpha_j)\in\Sigma_j^{D_j}\).
          \item \(\prover,\verifier,\mathcal F_{\mathrm{MsgOracle}}\): run
                \(\mathcal F_{\mathrm{MsgOracle}}.\mathsf{Commit}(\pi_j)\),
                producing \(O_j\).
        \end{enumerate}
  \item \(\verifier\): choose a finite input-opening set \(Q_y\subseteq D\) and
        finite prover-message opening sets \(Q_j\subseteq D_j\) for
        \(j\in\{0,\ldots,T\}\).
  \item \(\verifier\): for each \(a\in Q_y\), call
        \(\mathcal F_{\mathrm{MsgOracle}}.\mathsf{Open}(O_y,a)\).
  \item \(\verifier\): for each \(j\in\{0,\ldots,T\}\) and each
        \(a\in Q_j\), call
        \(\mathcal F_{\mathrm{MsgOracle}}.\mathsf{Open}(O_j,a)\).
  \item \(\verifier\): output \(0\) or \(1\) from \(x\), the transcript,
        public challenges, and opened values.
  \item \(\mathsf{Completeness}\): if \(y\in S_{\mathrm{valid}}(x)\), the
        honest prover is accepted.
  \item \(\mathsf{Proximity\ soundness}\): if
        \(\Delta(y,S_{\mathrm{valid}}(x))>\delta\), every malicious prover is
        accepted with probability at most \(\varepsilon\).
  \item \(\mathsf{Gap}\): if \(0<\Delta(y,S_{\mathrm{valid}}(x))\le\delta\),
        the template makes no accept/reject promise.
\end{enumerate}
\end{protocolbox}

The important negative statement is as useful as the positive one: an IOPP does
not usually prove that the input string is exactly valid.  It proves that a
string far from the valid set is rejected.  For code-based protocols this means
``close to a codeword,'' not ``equal to a codeword.''

\subsection{Holographic Access Use Cases}

Holographic access is the same handle interface in each use case:
\[
  O_Y\quad\text{opens selected coordinates of}\quad Y\in\Sigma^D.
\]
What changes is why \(Y\) is present in the proof.

For an indexed family of relations~[COS20], \(Y=\mathsf{idx}(i)\) is a
preprocessed description of a large public object \(i\).  The symbol \(i\) is
not a coordinate index; it is the relation description, such as a circuit,
constraint system, fixed table, or wiring description.  The relation can be
written either as \(R(i,x,w)\) or as \(R_i(x,w)\).  The verifier does not read
all of \(i\).  Instead, an indexer preprocesses \(i\) into an openable string
\(\mathsf{idx}(i)\in\Gamma^{D_{\mathrm{idx}}}\), fixes that string behind a
message-oracle handle \(O_i\), and gives the handle to the verifier.

For a proximity proof, \(Y\) is the long statement string itself.  The verifier
has a handle \(O_Y\), and the protocol declares a valid set
\(S_{\mathrm{valid}}(x)\subseteq\Sigma^D\) for that string.  The proof should
reject if the literal string behind \(O_Y\) is far from
\(S_{\mathrm{valid}}(x)\).  Later code-proximity sections instantiate this
generic \(Y\) with their concrete oracle strings.

For prover messages, \(Y=\pi_j\) is a proof oracle string sent in round \(j\).
It is fixed before later verifier challenges or query choices can depend on
it.  This is the ordinary IOP use of the same coordinate-opening handle.

Thus there is no separate index-oracle API in this paper.  Index handles,
long-statement handles, prover-message handles, Merkle handles, and virtual
handles all implement the same coordinate-opening interface.  The surrounding
relation tells us whether the handle is static relation data, statement data,
or proof data.

\begin{protocolbox}{Template: Holographic IOP}{fig:template-holographic-iop}
Let \(R\subseteq\mathcal I\times\mathcal X\times\mathcal W\), let
\(\mathsf{idx}:\mathcal I\to\Gamma^{D_{\mathrm{idx}}}\) be the indexer's
preprocessing map, and let \(\varepsilon\in[0,1]\) be the declared soundness
error.  The handle \(O_i\) opens the fixed string \(\mathsf{idx}(i)\).

\smallskip
\noindent\(\underline{\mathsf{HolographicIOP}[R,\mathsf{idx},\varepsilon](\verifier:(x,O_i),\ \prover:(i,x,w))\to\{0,1\}}\)

\smallskip
\begin{enumerate}[leftmargin=0.24in,label=\textbf{G\arabic*.},itemsep=0.08em,topsep=0.1em,parsep=0pt]
  \item \(O_i\): fixed message-oracle handle implementing
        \(\mathcal F_{\mathrm{MsgOracle}}.\mathsf{Open}\).
  \item \(\prover\): compute
        \(\pi_0=f_0(i,x,w)\in\Sigma_0^{D_0}\).
  \item \(\prover,\verifier,\mathcal F_{\mathrm{MsgOracle}}\): run
        \(\mathcal F_{\mathrm{MsgOracle}}.\mathsf{Commit}(\pi_0)\),
        producing \(O_0\).
  \item For each round \(j\in[T]\):
        \begin{enumerate}[leftmargin=0.2in,label=\textbf{G\arabic{enumi}.\arabic*.},itemsep=0.05em,topsep=0.05em,parsep=0pt]
          \item \(\verifier\): sample public challenge \(\alpha_j\) from its
                declared challenge domain.
          \item \(\prover\): compute
                \(\pi_j=f_j(i,x,w,\alpha_1,\ldots,\alpha_j)\in\Sigma_j^{D_j}\).
          \item \(\prover,\verifier,\mathcal F_{\mathrm{MsgOracle}}\): run
                \(\mathcal F_{\mathrm{MsgOracle}}.\mathsf{Commit}(\pi_j)\),
                producing \(O_j\).
        \end{enumerate}
  \item \(\verifier\): choose a finite index-opening set
        \(Q_i\subseteq D_{\mathrm{idx}}\) and finite prover-message opening sets
        \(Q_j\subseteq D_j\) for \(j\in\{0,\ldots,T\}\).
  \item \(\verifier\): for each \(a\in Q_i\), call
        \(\mathcal F_{\mathrm{MsgOracle}}.\mathsf{Open}(O_i,a)\).
  \item \(\verifier\): for each \(j\in\{0,\ldots,T\}\) and each
        \(a\in Q_j\), call
        \(\mathcal F_{\mathrm{MsgOracle}}.\mathsf{Open}(O_j,a)\).
  \item \(\verifier\): output \(0\) or \(1\) from \(x\), the transcript,
        public challenges, and opened values.
  \item \(\mathsf{Completeness}\): if \((i,x,w)\in R\), the honest prover is
        accepted.
  \item \(\mathsf{Soundness}\): if there is no \(w\in\mathcal W\) such that
        \((i,x,w)\in R\), every malicious prover is accepted with probability
        at most \(\varepsilon\).
\end{enumerate}
\end{protocolbox}

The indexed template in Figure~\ref{fig:template-holographic-iop} is one
standard use case.  It has ordinary exact soundness for the indexed relation
\(R(i,x,w)\), assuming \(O_i\) is the fixed handle for the declared
preprocessing \(\mathsf{idx}(i)\).  The proximity use case is different: there
the fixed handle opens the long statement string whose distance from a valid set
is being tested.

\begin{center}
\small
\begin{tabular}{p{0.17\linewidth}|p{0.18\linewidth}|p{0.29\linewidth}|p{0.25\linewidth}}
\(\mathsf{Use}\) & \(\mathsf{Verifier\ input}\) & \(\mathsf{Handle\ opens}\) & \(\mathsf{Soundness\ promise}\) \\
\hline
IOP prover message & \(x\) plus transcript handles & \(\pi_j\), proof data & membership in \(R\) \\
IOPP statement string & \((x,O_y)\) & \(y\in\Sigma^D\), tested statement data & proximity gap for \(S_{\mathrm{valid}}(x)\) \\
Indexed relation data & \((x,O_i)\) & \(\mathsf{idx}(i)\), preprocessed relation data & membership in \(R(i,\cdot,\cdot)\)
\end{tabular}
\end{center}

In later Blaze sections, large structured objects such as permutations may be
treated as indexed relation data when they are not read directly by the
verifier.  Encoded words and virtual folded words may be long statement
strings.  Backend proof oracles may be prover-message strings.  Each later
section must name the concrete object and say which role it is playing.

\subsection{MLPCS Handles}

The protocols in Blaze are multilinear-polynomial IOPs before they are compiled
to proximity proofs and cryptographic commitments~[Blaze24].  Blaze calls this
model an MLIOP.  In this paper we use the name MLPCS for the same ideal
multilinear evaluation interface, because the operation we expose is a
commit-and-evaluate interface for a multilinear polynomial.  This is a naming
choice, not a protocol change.

In this template, the prover does not merely send field elements as messages.
The prover commits multilinear oracle tensors to an ideal functionality, and
the verifier later evaluates their multilinear extensions at field points.  The
native MLPCS model is exact: a committed tensor \(h\in\F^{\B^d}\) determines one
hatted evaluation handle \(\widehat O_h:\F^d\to\F\).  It is not yet a
proximity model or a concrete cryptographic PCS.

We use tensor notation for the state of such an oracle.  A tensor
\(h\in\F^{\B^d}\) determines one multilinear extension
\[
  \MLE(h):\F^d\to\F.
\]
The MLE oracle handle returned by the functionality is typed
\[
  \widehat O_h:\F^d\to\F,
  \qquad
  \widehat O_h(r)=\MLE(h)(r).
\]
The hat is part of the notation: \(\widehat O_h\) is MLE-evaluation access,
not coordinate-opening access.

\begin{protocolbox}{Functionality: \(\mathcal F_{\mathrm{MLPCS}}\) Oracle}{fig:func-mlpcs-oracle}
Let \(d\ge 0\).

\smallskip
\noindent\(\underline{\mathcal F_{\mathrm{MLPCS}}.\mathsf{Commit}(\prover:h\in\F^{\B^d},\ \verifier:d)\to \widehat O_h:\F^d\to\F}\)

\smallskip
\begin{enumerate}[leftmargin=0.24in,label=\textbf{L\arabic*.},itemsep=0.08em,topsep=0.1em,parsep=0pt]
  \item \(\mathcal F_{\mathrm{MLPCS}}\): choose a fresh MLE handle
        \(\widehat O_h\) and store the committed tensor
        \(h\in\F^{\B^d}\) under that handle.
  \item \(\mathcal F_{\mathrm{MLPCS}}\): define
        \[
          \widehat O_h(r)
          :=
          \sum_{b\in\B^d} h[b]\,\chi_b(r)
          \qquad(r\in\F^d).
        \]
  \item \(\mathcal F_{\mathrm{MLPCS}}\): return \(\widehat O_h\).
\end{enumerate}

\smallskip
\noindent\(\underline{\mathcal F_{\mathrm{MLPCS}}.\mathsf{Eval}(\verifier:(\widehat O_h:\F^d\to\F,\ r\in\F^d))\to(\verifier:\widehat O_h(r)\in\F)}\)

\smallskip
\begin{enumerate}[leftmargin=0.24in,label=\textbf{E\arabic*.},itemsep=0.08em,topsep=0.1em,parsep=0pt]
  \item \(\mathcal F_{\mathrm{MLPCS}}\): return
        \((\verifier:\widehat O_h(r))\).
\end{enumerate}
\end{protocolbox}

This is not a coordinate-opening interface.  The prover commits the whole
native MLPCS tensor \(h\), and the output \(\widehat O_h\) is a
field-evaluation oracle.  If \(r\in\B^d\), then
\(\widehat O_h(r)=h[r]\), so a Boolean evaluation happens to return a tensor
entry.  But for general \(r\in\F^d\), the functionality returns an interpolated
MLE value.  This distinction is exactly why an MLPCS is not already a
Merkle-based PCS.

\begin{protocolbox}{Template: MLPCS}{fig:template-mlpcs}
Let \(R\subseteq\mathcal I\times\mathcal X\times\mathcal Y\times\mathcal W\),
where \(x\in\mathcal X\) is read directly and \(y\in\mathcal Y\) is a named
long statement tensor.

\smallskip
\noindent\(\underline{\mathsf{MLPCS}[R](\verifier:(i,x,\widehat O_y),\ \prover:(i,x,y,w))\to\{0,1\}}\)

\smallskip
\begin{enumerate}[leftmargin=0.24in,label=\textbf{A\arabic*.},itemsep=0.08em,topsep=0.1em,parsep=0pt]
  \item \(\widehat O_y:\F^d\to\F\): implements MLE access to the declared
        long statement tensor \(y\in\F^{\B^d}\).
  \item \(\prover,\verifier,\mathcal F_{\mathrm{MLPCS}}\): commit one
        or more prover-chosen tensors to \(\mathcal F_{\mathrm{MLPCS}}\),
        producing hatted MLE handles
        \(\widehat O_{h_1},\ldots,\widehat O_{h_s}\).
  \item \(\verifier\): sample public challenges.
  \item \(\prover,\verifier\): repeat MLPCS rounds as specified by the
        protocol.
  \item \(\verifier\): evaluate MLPCS handles through
        \(\mathcal F_{\mathrm{MLPCS}}.\mathsf{Eval}\) at specified field
        points.
  \item \(\verifier\): accept or reject from the transcript, challenges, and
        evaluated values.
\end{enumerate}
\end{protocolbox}

The word PIOP is the more general phrase ``polynomial interactive oracle
proof,'' as in HyperPlonk~[CBBZ23].  Blaze uses the term MLIOP for the
multilinear specialization.  Our MLPCS name refers to that same idealized
object: every prover polynomial oracle is implemented by an ideal functionality
storing a prover-committed tensor
\(h\in\F^{\B^d}\) and answering evaluations of \(\MLE(h)\).  Equivalently, it
idealizes the map
\[
  h\in\F^{\B^d}
  \quad\leadsto\quad
  \widehat O_h:\F^d\to\F.
\]

\subsection{Holographic MLPCS MLE Handles}

The native functionality above allows the prover to commit any tensor \(h\).
This is faithful to the native MLPCS model, but it is not enough to describe
the systematic bridge in Blaze's MLIOP-to-IOPP compiler.  In Blaze's literal
presentation, the long statement string \(y\in\F^{\B^d}\) is already available
as the input oracle, and the compiler uses a systematic backend code
\(C_{\mathrm{sys}}(h)=h\Vert\widetilde h\): the prover supplies the
non-systematic part while the systematic/message part is the original input
oracle.  Our handle API names the same binding as a separate functionality.

This is different from \(\EncMLE[C]\).  In \(\EncMLE[C]\), the verifier gives a
handle to a full encoded string and asks whether that string is close to
\(C(m)\) while \(\MLE(m)(z)=v\).  In \(\mathcal F_{\mathrm{HMLPCS}}^\delta\),
the verifier gives a handle \(O_y\) to the message-domain string itself,
\(y\in\F^{\B^d}\).  The prover supplies an effective message tensor
\(h\in\F^{\B^d}\), and the functionality returns evaluations of
\(\MLE(h)\) only when \(h\) is close to the concrete string behind \(O_y\).
Thus HMLPCS is the message-domain interface; a systematic encoded-MLE backend
is one way to realize it.

The prover commits an effective tensor
\(h\in\F^{\B^d}\).  The commit call always creates an attempted MLE handle
\(\widehat O_h\).  The proximity decision is enforced when that handle is
evaluated: if \(h\) is compatible with the string behind \(O_y\), evaluation
returns \(\MLE(h)(r)\); if it is too far, evaluation returns \(\bot\).
For a message-oracle handle \(O_y\) implementing
\(\mathcal F_{\mathrm{MsgOracle}}.\mathsf{Open}\), define
\[
  \Delta(h,O_y)
  =
  \frac{
    |\{b\in\B^d:
      h[b]\ne \mathcal F_{\mathrm{MsgOracle}}.\mathsf{Open}(O_y,b)\}|
  }{2^d}.
\]

\begin{protocolbox}{Functionality: \(\mathcal F_{\mathrm{HMLPCS}}^\delta\) Oracle}{fig:func-hmlpcs-oracle}
\noindent Public parameters: a field \(\F\), a dimension \(d\ge 0\), and a
frontend proximity threshold \(\delta\in[0,1]\).

\smallskip
\noindent\(\underline{\mathcal F_{\mathrm{HMLPCS}}^\delta.\mathsf{Commit}(\prover:h\in\F^{\B^d},\ \verifier:O_y:\B^d\to\F\cup\{\bot\})\to \widehat O_h}\)

\smallskip
\begin{enumerate}[leftmargin=0.24in,label=\textbf{H\arabic*.},itemsep=0.08em,topsep=0.1em,parsep=0pt]
  \item \(\mathcal F_{\mathrm{HMLPCS}}^\delta\): choose a fresh MLE handle
        \(\widehat O_h\).  Its evaluation interface has type
        \(\widehat O_h:\F^d\to\F\cup\{\bot\}\).
  \item \(\mathcal F_{\mathrm{HMLPCS}}^\delta\): store the committed tensor
        \(h\in\F^{\B^d}\), the coordinate handle \(O_y\), and the proximity
        threshold \(\delta\) under \(\widehat O_h\).
  \item \(\mathcal F_{\mathrm{HMLPCS}}^\delta\): return \(\widehat O_h\).
\end{enumerate}

\smallskip
\noindent\(\underline{\mathcal F_{\mathrm{HMLPCS}}^\delta.\mathsf{Eval}(\verifier:(\widehat O_h,\ r\in\F^d))\to(\verifier:\F\cup\{\bot\})}\)

\smallskip
\begin{enumerate}[leftmargin=0.24in,label=\textbf{E\arabic*.},itemsep=0.08em,topsep=0.1em,parsep=0pt]
  \item \(\mathcal F_{\mathrm{HMLPCS}}^\delta\): if \(\widehat O_h\) was not
        previously returned by a commit call, return \(\bot\).
  \item \(\mathcal F_{\mathrm{HMLPCS}}^\delta\): retrieve
        \(h,O_y,\delta\) stored under \(\widehat O_h\), read the concrete tensor
        \(y\in\F^{\B^d}\) through \(O_y\), and compute
        \(\Delta_h\gets\Delta(h,O_y)\).
  \item If \(\Delta_h>\delta\), return \(\bot\).
  \item If \(\Delta_h=0\), return
        \[
          \MLE(h)(r)=\sum_{b\in\B^d}h[b]\chi_b(r).
        \]
  \item If \(0<\Delta_h\le\delta\), let the corrupt source, or equivalently
        the simulator for that source, choose whether to return \(\bot\) or to
        return \(\MLE(h)(r)\).
\end{enumerate}
\end{protocolbox}

The returned hatted handle evaluates the effective committed tensor \(h\), not
the coordinate handle \(O_y\).  The proximity test is in the message domain:
\(\Delta(h,O_y)\) compares \(h[b]\) directly with the value opened from
\(O_y\) at the same \(b\in\B^d\).  No frontend code \(C\) is part of the
HMLPCS functionality.

A concrete realization can use a systematic encoded-MLE opening.  Let
\[
  C_{\mathrm{sys}}:\F^{\B^d}\to\F^{D_{\mathrm{msg}}\sqcup D_{\mathrm{par}}}
\]
be systematic, with \(D_{\mathrm{msg}}=\B^d\) and
\[
  C_{\mathrm{sys}}(h)|_{D_{\mathrm{msg}}}=h.
\]
The verifier's existing coordinate handle \(O_y\) supplies the systematic
coordinates.  The prover supplies a parity handle \(O_{\mathrm{par}}\) for the
remaining coordinates.  Together they define a virtual encoded handle
\[
  O_{\mathrm{enc}}(a)=
  \begin{cases}
    O_y(a) & a\in D_{\mathrm{msg}},\\
    O_{\mathrm{par}}(a) & a\in D_{\mathrm{par}}.
  \end{cases}
\]
An HMLPCS evaluation claim \(\MLE(h)(r)=v\) is then implemented by running
\[
  \mathcal F_{\EncMLE[C_{\mathrm{sys}}]}^{\delta_{\mathrm{enc}}}
    .\mathsf{Open}(O_{\mathrm{enc}},r,v).
\]
The proximity threshold must be chosen for the message-domain guarantee.  If
\[
  \alpha=
  \frac{|D_{\mathrm{msg}}|}
       {|D_{\mathrm{msg}}\sqcup D_{\mathrm{par}}|}
\]
is the systematic fraction of the encoded word, then full-codeword proximity
\(\delta_{\mathrm{enc}}\) only implies message-domain proximity
\(\delta_{\mathrm{enc}}/\alpha\).  To realize
\(\mathcal F_{\mathrm{HMLPCS}}^\delta\), choose
\[
  \delta_{\mathrm{enc}}\le \alpha\delta.
\]
This is the concrete reason the backend proximity parameter for the
holographic endpoint may need to be smaller than the proximity parameter one
would use for an ordinary encoded-MLE opening.

\subsection{Message-Oracle Handles}

The oracle functionalities above are the ideal parties used by these templates.
They give the verifier opening or evaluation access to committed objects.
A compiled protocol replaces those ideal functionalities with cryptographic
commitments and openings, as in the Merkle/Fiat--Shamir compilations used by
BaseFold, Blaze, and Fractal~[BaseFold23,Blaze24,COS20].  In this paper, the
compiled coordinate-opening interface is still the message-oracle interface from
Figure~\ref{fig:func-msg-oracle}.

For clarity, we will call any implementation handle for that interface a
\(\mathsf{MsgOracleHandle}\).  It implements coordinate openings:
\[
  \mathsf{Open}(C_x,a)\to x[a]\quad\text{or}\quad\bot.
\]
It does not implement arbitrary MLE evaluation.  If a later protocol needs to
authenticate
\[
  \MLE(h)(r)=v,
\]
it must reduce that claim to coordinate openings and algebraic checks, or call
a previously defined authentication backend.

The following table is the boundary discipline used in the rest of the paper.
Each row names the interface, not an implementation detail.

\begin{center}
\small
\begin{tabular}{p{0.19\linewidth}p{0.27\linewidth}p{0.42\linewidth}}
\textbf{Interface} & \textbf{Operation} & \textbf{Meaning} \\
\hline
\(\mathsf{MsgOracleHandle}\) &
\(\mathsf{Open}(C,a)\to x[a]\) or \(\bot\) &
Coordinate access to a fixed tensor/string.  Merkle, composite, and virtual
handles all implement this same interface. \\
\(\mathcal F_{\mathrm{MLPCS}}\) handle &
\(\mathsf{Commit}(h)\to\widehat O_h\), then
\(\mathsf{Eval}(\widehat O_h,r)\to\widehat O_h(r)\) &
Native MLPCS access to a prover-committed tensor
\(h\in\F^{\B^d}\).  This exists only in the native model before compilation. \\
\(\mathcal F_{\mathrm{HMLPCS}}^\delta\) handle &
\(\mathsf{Commit}(h,O_y)\to\widehat O_h\), then
\(\mathsf{Eval}(\widehat O_h,r)\to\MLE(h)(r)\) or \(\bot\) &
Proximity-bound MLE access: the effective tensor \(h\) must be close to the
concrete string behind the fixed coordinate handle \(O_y\). \\
\(\mathcal F_{\EncMLE[C]}^\delta\) opening &
\(\mathsf{Open}(O_y,z,v)\to\{0,1\}\) &
Encoded multilinear-evaluation opening against an already-fixed coordinate
oracle.  Adding a commit step \(h\mapsto O_{C(h)}\) gives the ordinary MLPCS
interface. \\
\text{PCS implementation} &
\(\mathsf{Commit}\), opening proofs, verifier checks &
Concrete cryptographic realization of the MLPCS or HMLPCS evaluation-handle
interface, typically built from encoded-MLE openings.  PCS is not a separate
ideal functionality in this paper. \\
\end{tabular}
\end{center}

We will keep the interfaces separated.  The message-oracle functionality has
\(\mathsf{Commit}\) and \(\mathsf{Open}\): a source commits a full string and
the functionality returns a handle.  A concrete protocol may expose a method
named \(\mathsf{Commit}\), such as
\(\mathsf{RowFold.Commit}\).  In this paper that method is
only a constructor for a \(\mathsf{MsgOracleHandle}\).  After construction, the
coordinate interface is still exactly
\(\mathsf{Open}(C,a)\to x[a]\) or \(\bot\).

\subsection{Virtual Handles}

A handle can be virtual without changing its interface.  Virtual means that the
answer at an address is computed from other handles instead of being stored as
one literal tensor.  There is no new proof model or new trusted functionality
here.  A virtual handle is just an implementation of an existing open interface.

We use one verb for this interface: \(\mathsf{Open}(O,a)\).  For an ideal
message-oracle handle \(O\), this means
\(\mathcal F_{\mathrm{MsgOracle}}.\mathsf{Open}(O,a)\).  For a compiled Merkle
or virtual handle, it means the authenticated opening operation.  In either
case the return type is a value or \(\bot\).

\begin{protocolbox}{Definition: Virtual Oracle Opening}{fig:def-virtual-oracle-opening}
Let \(O_{\mathrm{virt}}\) be a virtual handle with address set \(D\) and value
space \(\Sigma\).  Its implementation is fixed by backend handles
\(O_1,\ldots,O_m\).  For each \(a\in D\), it specifies a finite opening plan
\[
  Q(a)=((O^{(1)},a_1),\ldots,(O^{(s)},a_s)),
\]
where each \(O^{(\nu)}\in\{O_1,\ldots,O_m\}\) and \(a_\nu\) is an address for
that backend handle.  It also specifies a deterministic transformation
\[
  \Phi_a:\Sigma_1'\times\cdots\times\Sigma_s'\to\Sigma,
\]
where \(\Sigma_\nu'\) is the value space of \(O^{(\nu)}\).

\smallskip
\noindent\(\underline{\mathsf{VirtualOracle.Open}(O_{\mathrm{virt}},a\in D)\to\Sigma\cup\{\bot\}}\)

\smallskip
\begin{enumerate}[leftmargin=0.24in,label=\textbf{V\arabic*.},itemsep=0.08em,topsep=0.1em,parsep=0pt]
  \item Let \(Q(a)=((O^{(1)},a_1),\ldots,(O^{(s)},a_s))\).
  \item For each \(\nu\in\{1,\ldots,s\}\), set
        \(v_\nu\gets\mathsf{Open}(O^{(\nu)},a_\nu)\).
  \item If any \(v_\nu=\bot\), return \(\bot\).
  \item Return \(\Phi_a(v_1,\ldots,v_s)\).
\end{enumerate}
\end{protocolbox}

The most common virtual handle in this paper is a same-address linear
combination, which we call \(\mathsf{RowFold}\).  Let
\[
  O_0,\ldots,O_{q-1}
\]
be handles with the same address set \(D\) and scalar value space \(\F\).  Let
\[
  \lambda_0,\ldots,\lambda_{q-1}\in\F
\]
be public coefficients.  The constructor
\[
  O_{\mathrm{fold}}
  \gets
  \mathsf{RowFold.Commit}\bigl((O_j,\lambda_j)_{j=0}^{q-1}\bigr)
\]
returns a virtual handle over the same address set \(D\).  Its opening rule is
\[
  \mathsf{Open}(O_{\mathrm{fold}},a)
  =
  \sum_{j=0}^{q-1}
    \lambda_j\,\mathsf{Open}(O_j,a),
  \qquad a\in D,
\]
with the convention that the virtual opening returns \(\bot\) if any required
backend opening returns \(\bot\).  The prover does not commit a fresh tensor
for \(O_{\mathrm{fold}}\); the folded handle is verifier-derived from existing
handles and public coefficients.

The same definition applies to any layer whose handle interface has a
same-address read operation.  For message-oracle handles, the read operation is
\(\mathsf{Open}\).  For an MLE handle
\(\widehat O_h:\F^d\to\F\), the read operation is evaluation at a field point.
Thus if \(\widehat O_0,\ldots,\widehat O_{q-1}\) have the same domain \(\F^d\),
then
\[
  \widehat O_{\mathrm{fold}}
  \gets
  \mathsf{RowFold.Commit}\bigl((\widehat O_j,\lambda_j)_{j=0}^{q-1}\bigr)
\]
means the virtual MLE handle satisfying
\[
  \widehat O_{\mathrm{fold}}(r)
  =
  \sum_{j=0}^{q-1}\lambda_j\,\widehat O_j(r),
  \qquad r\in\F^d.
\]

The canonical Blaze example is a row combination.  Suppose
\[
  c\in\F^{\B^\ell\times\B^n}
  \qquad\text{and}\qquad
  \rho\in\F^{\B^\ell}.
\]
The logical folded tensor \(c_{\mathrm{combo}}\in\F^{\B^n}\) is
\[
  c_{\mathrm{combo}}[j]
  =
  \sum_{h\in\B^\ell}\rho[h]\,c[h,j].
\]
Equivalently, for each \(h\in\B^\ell\), view the \(h\)-th row as a handle
\[
  O_h(j)=c[h,j],\qquad j\in\B^n.
\]
Then
\[
  O_{\mathrm{virt}}
  =
  \mathsf{RowFold.Commit}\bigl((O_h,\rho[h])_{h\in\B^\ell}\bigr)
\]
is the virtual handle for \(c_{\mathrm{combo}}\).  If the implementation stores
the matrix as one packed column handle \(O_c\), the row handles \(O_h\) are
implemented by opening \(O_c\) at the same column and projecting lane \(h\).
In the notation of Figure~\ref{fig:def-virtual-oracle-opening}, the opening plan
for address \(j\) is
\[
  Q(j)=((O_c,(h,j)))_{h\in\B^\ell},
\]
and the transformation is
\[
  \Phi_j((v_h)_{h\in\B^\ell})
  =
  \sum_{h\in\B^\ell}\rho[h]v_h.
\]
Thus \(\mathsf{VirtualOracle.Open}(O_{\mathrm{virt}},j)\) opens the column
\(c[\ast,j]\), computes the displayed linear combination, and returns
\(c_{\mathrm{combo}}[j]\).  If the backend handle \(O_c\) is an ideal
message-oracle handle, these backend calls are ideal
\(\mathcal F_{\mathrm{MsgOracle}}.\mathsf{Open}\) calls.  If \(O_c\) is a
Merkle or virtual implementation handle, these backend calls are authenticated
openings implementing the same interface.

Virtuality is therefore not a new trust assumption.  It is an implementation of
an already-defined handle interface.

\subsection{Layer Checklist}

The rest of the paper should be read with this checklist in mind.

\begin{enumerate}[label=\textbf{H\arabic*.},leftmargin=0.28in]
  \item If a protocol says ``oracle string,'' it should mean a string stored by
        \(\mathcal F_{\mathrm{MsgOracle}}\), or a handle implementing the same
        open interface.
  \item If a protocol says ``multilinear oracle,'' it should mean a hatted MLE
        handle \(\widehat O_h:\F^d\to\F\) produced by
        \(\mathcal F_{\mathrm{MLPCS}}.\mathsf{Commit}\) for a native
        prover-committed tensor, or by
        \(\mathcal F_{\mathrm{HMLPCS}}^\delta.\mathsf{Commit}\) when the
        tensor is proximity-checked against a verifier coordinate handle.
  \item If a protocol says ``commitment,'' it should mean a concrete
        \(\mathsf{MsgOracleHandle}\) implementing
        Figure~\ref{fig:func-msg-oracle}.
  \item If a protocol exposes a method named \(\mathsf{Commit}\), it must state
        which handle it returns and which \(\mathsf{Open}\) interface that
        handle implements.
  \item If a protocol says ``virtual,'' it must still name the interface being
        implemented and the devirtualization rule.
  \item If a protocol moves from MLPCS to IOPP or from IOPP to PCS, it must say
        which ideal access has been replaced and what soundness theorem is being
        consumed.
\end{enumerate}
